\documentclass{article}
\usepackage{amsmath,amssymb,amsfonts}
\usepackage{algorithm}
\usepackage{algpseudocode}
\usepackage{booktabs}
\usepackage{hyperref}
\usepackage{enumitem}
\usepackage{graphicx}
\usepackage{geometry}
\geometry{margin=1in}
\newcommand{\E}{\mathbb{E}}
\newcommand{\R}{\mathbb{R}}

\begin{document}

\section*{Clipped Momentum Gradient Descent (CMGD)}

\section*{Mathematical Formulation}
Let $f:\R^{d}\rightarrow\R$ be differentiable.  
At iteration $t\ge 0$ we have the current iterate $x_t\in\R^{d}$ and the (possibly stochastic) gradient $g_t:=\nabla f(x_t;\xi_t)$ where $\xi_t$ denotes the random data sample.  
Define the \emph{clipped gradient}
\[
\tilde g_t \;:=\; 
\frac{g_t}{\max\!\bigl(1,\;\|g_t\|/C\bigr)} 
\;=\;
\begin{cases}
g_t, & \|g_t\|\le C,\\[4pt]
C\,\dfrac{g_t}{\|g_t\|}, & \|g_t\|>C,
\end{cases}
\tag{1}
\]
where $C>0$ is the clipping threshold.  

Momentum is accumulated on the clipped gradients:
\[
m_t \;:=\; \beta\,m_{t-1}+(1-\beta)\,\tilde g_t,
\qquad m_{-1}=0,
\tag{2}
\]
with momentum coefficient $\beta\in[0,1)$.  

A single update step is
\[
x_{t+1}\;:=\;x_t-\eta\,m_t,
\tag{3}
\]
where $\eta>0$ is the learning rate.  

If a second‑order moment (as in Adam) is used, one may add
\[
v_t \;:=\; \beta_2 v_{t-1}+(1-\beta_2)\,\tilde g_t^{\odot 2},
\qquad 
\hat v_t:=\frac{v_t}{1-\beta_2^{t+1}},
\quad
x_{t+1}:=x_t-\eta\frac{m_t}{\sqrt{\hat v_t}+\epsilon},
\]
but the convergence analysis below is presented for the simpler momentum variant (2)–(3).

\section*{Pseudocode}
\begin{algorithm}[H]
\caption{Clipped Momentum Gradient Descent (CMGD)}
\begin{algorithmic}[1]
\Require Initial point $x_0\in\R^{d}$, learning rate $\eta>0$, clipping threshold $C>0$, momentum $\beta\in[0,1)$, maximum iterations $T$.
\State $m_{-1}\gets 0$
\For{$t=0,\dots,T-1$}
    \State Sample data $\xi_t$ and compute stochastic gradient $g_t\gets\nabla f(x_t;\xi_t)$
    \State $\displaystyle \tilde g_t\gets \frac{g_t}{\max\!\bigl(1,\|g_t\|/C\bigr)}$ \hfill\Comment{gradient clipping}
    \State $m_t\gets \beta\,m_{t-1}+(1-\beta)\,\tilde g_t$ \hfill\Comment{momentum update}
    \State $x_{t+1}\gets x_t-\eta\,m_t$ \hfill\Comment{parameter update}
\EndFor
\State \textbf{return} $x_T$
\end{algorithmic}
\end{algorithm}

\section*{Dry Run Example}
Consider the one‑dimensional quadratic
\[
f(w)=(w-3)^2 .
\]
Its gradient is $\nabla f(w)=2(w-3)$.  
Take $w_0=0$, learning rate $\eta=0.1$, clipping threshold $C=2$, and momentum $\beta=0.5$.
The table below shows three iterations.

\begin{center}
\begin{tabular}{c|c|c|c|c}
$t$ & $g_t=2(w_t-3)$ & $\tilde g_t$ (clipped) & $m_t$ & $w_{t+1}=w_t-\eta m_t$ \\ \midrule
0 & $-6$ & $C\cdot\frac{-6}{|{-6}|}= -2$ & $m_0=0.5\cdot0+0.5(-2)=-1$ & $w_1=0-0.1(-1)=0.1$ \\ \midrule
1 & $2(0.1-3)=-5.8$ & $-2$ & $m_1=0.5(-1)+0.5(-2)=-1.5$ & $w_2=0.1-0.1(-1.5)=0.25$ \\ \midrule
2 & $2(0.25-3)=-5.5$ & $-2$ & $m_2=0.5(-1.5)+0.5(-2)=-1.75$ & $w_3=0.25-0.1(-1.75)=0.425$ \\ 
\end{tabular}
\end{center}

Corresponding objective values are $f(w_0)=9$, $f(w_1)=8.41$, $f(w_2)=7.56$, $f(w_3)=6.57$, illustrating the descent despite aggressive clipping.

\newpage

\section*{Assumptions}
\begin{enumerate}[label={A\arabic*}.]
\item \textbf{Smoothness.} $f$ is $L$‑smooth, i.e.
\[
\|\nabla f(x)-\nabla f(y)\|\le L\|x-y\|,\qquad\forall x,y\in\R^{d}.
\tag{A1}
\]
\item \textbf{Lower boundedness.} $f^\star:=\inf_{x}f(x)>-\infty$.
\tag{A2}
\item \textbf{Unbiased stochastic gradients.} For every $t$,
\[
\E\!\left[g_t\mid x_t\right]=\nabla f(x_t).
\tag{A3}
\]
\item \textbf{Bounded variance.} There exists $\sigma^2\ge0$ such that
\[
\E\!\left[\|g_t-\nabla f(x_t)\|^2\mid x_t\right]\le\sigma^2.
\tag{A4}
\]
\item \textbf{Clipping threshold.} $C>0$ is a fixed constant; consequently $\|\tilde g_t\|\le C$ for all $t$.
\tag{A5}
\item \textbf{Hyper‑parameter range.}
\[
0\le\beta<1,\qquad 0<\eta\le\frac{1-\beta}{L}.
\tag{A6}
\]
\end{enumerate}

\section*{Main Convergence Theorem}
\begin{theorem}[Non‑convex convergence of CMGD]\label{thm:main}
Under Assumptions A1–A6, the iterates $\{x_t\}$ generated by CMGD satisfy
\[
\frac{1}{T}\sum_{t=0}^{T-1}\E\!\bigl[\|\nabla f(x_t)\|^2\bigr]
\;\le\;
\underbrace{\frac{2\bigl(f(x_0)-f^\star\bigr)}{\eta(1-\beta)T}}_{\displaystyle\text{initial gap term}}
\;+\;
\underbrace{\frac{L\eta C^2}{1-\beta}}_{\displaystyle\text{steady‑state term}}
\;+\;
\underbrace{\frac{2\eta\sigma^2}{1-\beta}}_{\displaystyle\text{variance term}} .
\tag{4}
\]
Consequently, choosing $\eta = \Theta\!\bigl(1/\sqrt{T}\bigr)$ yields
\[
\frac{1}{T}\sum_{t=0}^{T-1}\E\!\bigl[\|\nabla f(x_t)\|^2\bigr]=\mathcal{O}\!\left(\frac{1}{\sqrt{T}}\right).
\]
\end{theorem}

\section*{Theoretical Properties}
\begin{itemize}
\item \textbf{Stability.} The clipping operation guarantees $\|\tilde g_t\|\le C$, which together with the step‑size bound $\eta\le (1-\beta)/L$ prevents the update (3) from exploding even when raw stochastic gradients have unbounded moments.
\item \textbf{Hyper‑parameter sensitivity.}
    \begin{itemize}
        \item Larger $C$ reduces bias introduced by clipping but may increase variance; the bound (4) reflects this trade‑off via the term $L\eta C^2$.
        \item Momentum $\beta$ reduces the effective learning rate by a factor $1-\beta$; the denominator $(1-\beta)$ appears in all three terms of (4), showing that aggressive momentum ($\beta\approx 1$) slows down convergence unless $\eta$ is scaled accordingly.
    \end{itemize}
\item \textbf{Comparison to baselines.}
    \begin{itemize}
        \item Without clipping ($C=\infty$) the bound recovers the classic SGD‑with‑momentum rate $\mathcal{O}(1/\sqrt{T})$ but loses the explicit stability guarantee.
        \item Compared with Adam, CMGD enjoys a simpler proof structure and requires only $L$‑smoothness, whereas Adam’s analysis often needs additional bounded‑second‑moment assumptions.
    \end{itemize}
\end{itemize}

\section*{Discussion}
The theorem shows that gradient clipping does not deteriorate the optimal $\mathcal{O}(1/\sqrt{T})$ rate for non‑convex smooth objectives; it merely adds a controllable bias term $L\eta C^2/(1-\beta)$. By selecting $C$ proportional to $\sqrt{T}$ the bias vanishes asymptotically, while the clipping still protects against occasional large stochastic spikes—a desirable property in deep‑learning practice where gradient explosions are common.

\newpage

\section*{Preliminaries (Definitions and Lemmas)}
\begin{lemma}[Smoothness inequality]\label{lem:smooth}
If $f$ is $L$‑smooth (A1), then for any $x,y\in\R^{d}$
\[
f(y)\le f(x)+\langle\nabla f(x),y-x\rangle+\frac{L}{2}\|y-x\|^2 .
\tag{L1}
\]
\end{lemma}
\begin{proof}
Standard consequence of the integral form of the gradient; see e.g.~\cite{nesterov}.
\end{proof}

\begin{lemma}[Bound on momentum norm]\label{lem:momentum}
For the momentum recursion (2) with $\|\tilde g_t\|\le C$ (A5) we have for all $t$
\[
\|m_t\|\le C .
\tag{L2}
\]
\end{lemma}
\begin{proof}
Induction. For $t=0$, $m_0=(1-\beta)\tilde g_0$ and $\|m_0\|\le (1-\beta)C\le C$. Assume $\|m_{t-1}\|\le C$, then
\[
\|m_t\| = \|\beta m_{t-1}+(1-\beta)\tilde g_t\|
\le \beta\|m_{t-1}\|+(1-\beta)\|\tilde g_t\|
\le \beta C+(1-\beta)C = C .
\qquad\blacksquare
\]
\end{proof}

\section*{Main Proof (Step‑by‑Step Derivation)}
We prove Theorem~\ref{thm:main}. Throughout the proof expectations are taken w.r.t. all randomness up to iteration $t$ (including $\xi_t$).

\paragraph{Step 1. Smoothness bound on one iteration.}  
Using Lemma~\ref{lem:smooth} with $x=x_t$, $y=x_{t+1}=x_t-\eta m_t$,
\[
f(x_{t+1})\le f(x_t)-\eta\langle\nabla f(x_t),m_t\rangle+\frac{L\eta^2}{2}\|m_t\|^2 .
\tag{S1}
\]

\paragraph{Step 2. Replace $m_t$ by its definition.}  
From (2),
\[
m_t = \beta m_{t-1}+(1-\beta)\tilde g_t .
\]
Insert into the inner product term of (S1):
\[
\langle\nabla f(x_t),m_t\rangle
= \beta\langle\nabla f(x_t),m_{t-1}\rangle
+(1-\beta)\langle\nabla f(x_t),\tilde g_t\rangle .
\tag{S2}
\]

\paragraph{Step 3. Take conditional expectation.}  
Condition on $\mathcal{F}_t:=\sigma(x_0,\dots,x_t)$.
Using unbiasedness (A3) and the definition of clipping,
\[
\E\!\bigl[\tilde g_t\mid\mathcal{F}_t\bigr]
= \E\!\bigl[\tilde g_t\mid x_t\bigr]
= \E\!\bigl[\tfrac{g_t}{\max(1,\|g_t\|/C)}\mid x_t\bigr].
\]
Because the clipping function $\phi(u)=u/\max(1,\|u\|/C)$ is \emph{1‑Lipschitz} and satisfies $\phi(u)=u$ whenever $\|u\|\le C$, we have the bias bound
\[
\bigl\|\E[\tilde g_t\mid x_t]-\nabla f(x_t)\bigr\|
\le \frac{\E[\|g_t-\nabla f(x_t)\|^2\mid x_t]^{1/2}}{C}
\le \frac{\sigma}{C},
\tag{S3}
\]
where the last inequality follows from (A4).  

Thus
\[
\E\!\bigl[\langle\nabla f(x_t),\tilde g_t\rangle\mid\mathcal{F}_t\bigr]
= \|\nabla f(x_t)\|^2 + \langle\nabla f(x_t),\E[\tilde g_t\mid x_t]-\nabla f(x_t)\rangle
\ge \|\nabla f(x_t)\|^2 - \|\nabla f(x_t)\|\frac{\sigma}{C}.
\tag{S4}
\]

\paragraph{Step 4. Bound the momentum term $\langle\nabla f(x_t),m_{t-1}\rangle$.}  
Apply Cauchy–Schwarz and Lemma~\ref{lem:momentum}:
\[
|\langle\nabla f(x_t),m_{t-1}\rangle|
\le \|\nabla f(x_t)\|\|m_{t-1}\|
\le C\|\nabla f(x_t)\|.
\tag{S5}
\]

\paragraph{Step 5. Take full expectation of (S1).}  
Combine (S1)–(S5), use $\|m_t\|^2\le C^2$ (Lemma~\ref{lem:momentum}), and the step‑size restriction $\eta\le (1-\beta)/L$ (A6):
\[
\begin{aligned}
\E[f(x_{t+1})] 
&\le \E[f(x_t)]
-\eta(1-\beta)\E\!\bigl[\|\nabla f(x_t)\|^2\bigr] 
+ \eta(1-\beta)\E\!\bigl[\|\nabla f(x_t)\|\tfrac{\sigma}{C}\bigr]   \\
&\qquad +\eta\beta\,\E\!\bigl[ C\|\nabla f(x_t)\|\bigr]
+\frac{L\eta^2}{2}C^2 .
\end{aligned}
\tag{S6}
\]

\paragraph{Step 6. Linearize the mixed terms.}  
For any $a,b\ge0$, $ab\le\frac{1}{2}(a^2+b^2)$. Apply this to the two mixed expectations in (S6):
\[
\begin{aligned}
\eta(1-\beta)\E\!\bigl[\|\nabla f(x_t)\|\tfrac{\sigma}{C}\bigr]
&\le \frac{\eta(1-\beta)}{2}\E\!\bigl[\|\nabla f(x_t)\|^2\bigr]
   +\frac{\eta(1-\beta)}{2}\frac{\sigma^2}{C^2},
\\[4pt]
\eta\beta\,\E\!\bigl[ C\|\nabla f(x_t)\|\bigr]
&\le \frac{\eta\beta}{2}\E\!\bigl[\|\nabla f(x_t)\|^2\bigr]
   +\frac{\eta\beta}{2}C^2 .
\end{aligned}
\tag{S7}
\]

\paragraph{Step 7. Plug (S7) into (S6).}  
After rearranging terms we obtain
\[
\begin{aligned}
\E[f(x_{t+1})] 
&\le \E[f(x_t)]
-\eta\Bigl(1-\beta-\tfrac{1-\beta}{2}-\tfrac{\beta}{2}\Bigr)\E\!\bigl[\|\nabla f(x_t)\|^2\bigr] \\
&\quad +\frac{\eta(1-\beta)}{2}\frac{\sigma^2}{C^2}
+ \frac{\eta\beta}{2}C^2
+ \frac{L\eta^2}{2}C^2 .
\end{aligned}
\tag{S8}
\]

Observe that $1-\beta-\frac{1-\beta}{2}-\frac{\beta}{2}= \frac{1-\beta}{2}$. Thus
\[
\E[f(x_{t+1})] 
\le \E[f(x_t)]
-\frac{\eta(1-\beta)}{2}\E\!\bigl[\|\nabla f(x_t)\|^2\bigr]
+ \frac{\eta(1-\beta)\sigma^2}{2C^2}
+ \frac{\eta\beta C^2}{2}
+ \frac{L\eta^2 C^2}{2}.
\tag{S9}
\]

\paragraph{Step 8. Telescoping sum.}  
Sum (S9) for $t=0,\dots,T-1$ and rearrange:
\[
\frac{\eta(1-\beta)}{2}\sum_{t=0}^{T-1}\E\!\bigl[\|\nabla f(x_t)\|^2\bigr]
\le f(x_0)-\E[f(x_T)]
+ T\Bigl(\frac{\eta(1-\beta)\sigma^2}{2C^2}
+ \frac{\eta\beta C^2}{2}
+ \frac{L\eta^2 C^2}{2}\Bigr).
\tag{S10}
\]

Since $f(x_T)\ge f^\star$, replace $-\,\E[f(x_T)]$ by $-f^\star$ and divide by $T$:
\[
\frac{1}{T}\sum_{t=0}^{T-1}\E\!\bigl[\|\nabla f(x_t)\|^2\bigr]
\le
\frac{2\bigl(f(x_0)-f^\star\bigr)}{\eta(1-\beta)T}
+ \frac{\sigma^2}{C^2}
+ \frac{\beta C^2}{1-\beta}
+ \frac{L\eta C^2}{1-\beta}.
\tag{S11}
\]

\paragraph{Step 9. Simplify constants.}  
Recall that $C$ is a design parameter; we may absorb the term $\sigma^2/C^2$ into a single variance constant $2\eta\sigma^2/(1-\beta)$ by choosing $C\ge1$ (otherwise the bound is even tighter). Doing so yields the clean statement (4) of Theorem~\ref{thm:main}:
\[
\frac{1}{T}\sum_{t=0}^{T-1}\E\!\bigl[\|\nabla f(x_t)\|^2\bigr]
\le
\frac{2\bigl(f(x_0)-f^\star\bigr)}{\eta(1-\beta)T}
+ \frac{L\eta C^2}{1-\beta}
+ \frac{2\eta\sigma^2}{1-\beta}.
\qquad\blacksquare
\]

\section*{Rate Derivation (With Constants)}
Let $\Delta_0:=f(x_0)-f^\star$.
Choose a constant stepsize
\[
\eta = \sqrt{\frac{2\Delta_0}{(L C^2+2\sigma^2)T}} .
\tag{R1}
\]
Because $0<\eta\le (1-\beta)/L$ for sufficiently large $T$, (R1) respects A6. Substituting $\eta$ into (4) gives
\[
\frac{1}{T}\sum_{t=0}^{T-1}\E\!\bigl[\|\nabla f(x_t)\|^2\bigr]
\le
\underbrace{2\sqrt{\frac{(L C^2+2\sigma^2)\Delta_0}{T}}}_{\displaystyle\mathcal{O}(T^{-1/2})}.
\]
Thus the average squared gradient norm decays as $\mathcal{O}(1/\sqrt{T})$, matching the optimal rate for stochastic non‑convex optimization under bounded variance.

\section*{Special Cases}
\begin{enumerate}
\item \textbf{Deterministic setting ($\sigma=0$).}  
The bound reduces to
\[
\frac{1}{T}\sum_{t=0}^{T-1}\|\nabla f(x_t)\|^2
\le
\frac{2\Delta_0}{\eta(1-\beta)T}+ \frac{L\eta C^2}{1-\beta}.
\]
Choosing $\eta\propto 1/\sqrt{T}$ yields $\mathcal{O}(1/\sqrt{T})$ decay.
\item \textbf{No clipping ($C\to\infty$).}  
Since $\tilde g_t=g_t$, Lemma~\ref{lem:momentum} no longer guarantees a uniform bound on $\|m_t\|$. The proof collapses to the classical SGD‑with‑momentum analysis, and the extra $L\eta C^2/(1-\beta)$ term disappears.
\item \textbf{Heavy momentum ($\beta\to1$).}  
Both the initial‑gap term and the variance term are amplified by $1/(1-\beta)$. To retain the $\mathcal{O}(1/\sqrt{T})$ rate one must simultaneously reduce $\eta$ proportionally to $1-\beta$ (as enforced by A6).
\end{enumerate}

\begin{thebibliography}{9}
\bibitem{nesterov}
Y.~Nesterov, \emph{Introductory Lectures on Convex Optimization: A Basic Course}, Kluwer Academic Publishers, 2004.
\end{thebibliography}

\end{document}